% Chapter 1 -- Ordered Sets, 1.5-1.11
\rsection{Ordered Sets}

\begin{signpost}[Why this section is here]
Nothing in this section mentions $\R$, and on a first reading it feels like an
unmotivated detour into abstraction. It is not. Rudin has just shown that $\Q$
fails in a specific way, and he now needs language precise enough to \emph{name}
that failure --- because the name will become the defining property of $\R$ in
1.19. Definitions 1.5 to 1.7 are only scaffolding; 1.8 and 1.10 are the payload.
If you read one thing here, read Definition 1.10, and notice that it is exactly
what Example 1.9(a) showed $\Q$ does not have.
\end{signpost}

\ritem{1.5}{Definition}
Let $S$ be a set. An \emph{order} on $S$ is a relation, denoted by $<$, with the
following two properties:
\begin{enumerate}[label=(\roman*), leftmargin=2.2em, itemsep=2pt, topsep=4pt]
  \item If $x \in S$ and $y \in S$, then one and only one of the statements
        \[ x < y, \qquad x = y, \qquad y < x \]
        is true.\kn{Trichotomy. ``One and only one'' packs in two separate
        demands: at least one holds, so the order is \emph{total}; and at most
        one holds, so $<$ is irreflexive and antisymmetric. Drop totality and
        you get a partial order, where suprema can fail to exist for reasons
        having nothing to do with gaps.}
  \item If $x, y, z \in S$, if $x < y$ and $y < z$, then $x < z$.
\end{enumerate}
The statement ``$x < y$'' may be read as ``$x$ is less than $y$'' or ``$x$ is
smaller than $y$'' or ``$x$ precedes $y$''.

It is often convenient to write $y > x$ in place of $x < y$.

The notation $x \le y$ indicates that $x < y$ or $x = y$, without specifying
which of these two is to hold. In other words, $x \le y$ is the negation of
$x > y$.\kd{That last sentence is a small theorem, not a definition: ``$x<y$ or
$x=y$'' is equivalent to ``not $x>y$'' precisely because of trichotomy in (i).
In a partial order the two come apart, which is why $\le$ is normally taken as
the primitive notion there.}

\ritem{1.6}{Definition}
An \emph{ordered set} is a set $S$ in which an order is defined.

For example, $\Q$ is an ordered set if $r < s$ is defined to mean that $s - r$
is a positive rational number.

\ritem{1.7}{Definition}
Suppose $S$ is an ordered set, and $E \subset S$. If there exists a
$\beta \in S$ such that $x \le \beta$ for every $x \in E$, we say that $E$ is
\emph{bounded above}, and call $\beta$ an \emph{upper bound} of $E$.

Lower bounds are defined the same way (with $\ge$ in place of $\le$).

\ritem{1.8}{Definition}
Suppose $S$ is an ordered set, $E \subset S$, and $E$ is bounded above. Suppose
there exists an $\alpha \in S$ with the following properties:
\begin{enumerate}[label=(\roman*), leftmargin=2.2em, itemsep=2pt, topsep=4pt]
  \item $\alpha$ is an upper bound of $E$.
  \item If $\gamma < \alpha$ then $\gamma$ is not an upper bound of
        $E$.\kn{This is the clause you will actually use, and it is stated
        contrapositively. The working form: for every $\gamma < \alpha$ there
        is an $x \in E$ with $x > \gamma$. In $\R$, equivalently: for every
        $\varepsilon > 0$ there is $x \in E$ with $x > \alpha - \varepsilon$.
        Reach for that sentence every time a proof opens ``let
        $\alpha = \sup E$.''}
\end{enumerate}
Then $\alpha$ is called the \emph{least upper bound} of $E$ [that there is at
most one such $\alpha$ is clear from (ii)] or the \emph{supremum} of $E$, and we
write
\[ \alpha = \sup E. \]

The \emph{greatest lower bound}, or \emph{infimum}, of a set $E$ which is
bounded below is defined in the same manner: The statement
\[ \alpha = \inf E \]
means that $\alpha$ is a lower bound of $E$ and that no $\beta$ with
$\beta > \alpha$ is a lower bound of $E$.

\begin{unpack}[Why the least upper bound and not the largest element]
Because the largest element usually does not exist, and the supremum is the
repair. Take $E_1 = \{\, r \in \Q : r < 0 \,\}$. It has no largest element: for
any $r<0$ there is $r/2$ with $r < r/2 < 0$. But the set of its upper bounds ---
everything ${}\ge 0$ --- does have a smallest one, namely $0$.

The manoeuvre is to move the question from $E$, where the answer is often
``none,'' to the set of upper bounds of $E$, where it is often ``yes.'' Rudin's
bracketed aside is the payoff: uniqueness follows immediately from (ii), so
speaking of ``the'' supremum is legitimate.

\medskip
\begin{center}
\begin{tikzpicture}[x=1mm, y=1mm]
  % upper bounds, drawn above the axis so nothing sits on a line
  \draw[seg] (46,9) -- (86,9);
  \node[ptfill] at (46,9) {};
  \node[lbl, anchor=south] at (66,11) {upper bounds of $E$};
  % the axis
  \draw[axis] (-1,0) -- (92,0);
  % the set E
  \draw[seg] (6,0) -- (46,0);
  \node[ptopen] at (46,0) {};
  \node[lbl, anchor=north] at (26,-2) {the set $E$};
  % the seam
  \draw[guide] (46,-8) -- (46,15);
  \node[lbl, anchor=south] at (46,15) {$\alpha=\sup E$};
  % gamma
  \node[ptfill, minimum size=2.6pt] at (36,0) {};
  \node[mlbl, anchor=south] at (36,1.5) {$\gamma$};
  \draw[guide] (36,-8) -- (36,0);
  \draw[thin, {Stealth[length=3.5pt]}-{Stealth[length=3.5pt]}] (36,-8) -- (46,-8);
  \node[lbl, anchor=north] at (41,-9.5) {some $x\in E$ lies in here};
\end{tikzpicture}
\end{center}
\figcap{$\alpha$ is the seam. Everything to its right bounds $E$; nothing to its
left does.}

\smallskip
The picture also shows why clause (ii) is the useful one. Stepping back from
$\alpha$ by any amount lands you at a $\gamma$ that is \emph{not} an upper
bound, which is to say some element of $E$ sticks out past it. That is the
sentence you will use.
\end{unpack}

\ritem{1.9}{Examples}
\begin{enumerate}[label=(\alph*), leftmargin=2.2em, itemsep=5pt, topsep=4pt]
  \item Consider the sets $A$ and $B$ of Example 1.1 as subsets of the ordered
        set $\Q$. The set $A$ is bounded above. In fact, the upper bounds of $A$
        are exactly the members of $B$.\kd{Not \emph{exactly}, and Rudin is
        being loose. Every $\beta \in B$ bounds $A$, and no member of $A$ does.
        But one must also rule out a positive rational with $\beta^2 = 2$, and
        that is impossible only by Example 1.1. The claim holds, but it leans on
        the irrationality result proved a page earlier.} Since $B$ contains no
        smallest member, $A$ has no least upper bound in $\Q$.

        Similarly, $B$ is bounded below: the set of all lower bounds of $B$
        consists of $A$ and of all $r \in \Q$ with $r \le 0$. Since $A$ has no
        largest member, $B$ has no greatest lower bound in $\Q$.
  \item If $\alpha = \sup E$ exists, then $\alpha$ may or may not be a member of
        $E$. For instance, let $E_1$ be the set of all $r \in \Q$ with $r < 0$.
        Let $E_2$ be the set of all $r \in \Q$ with $r \le 0$. Then
        \[ \sup E_1 = \sup E_2 = 0, \]
        and $0 \notin E_1$, $0 \in E_2$.\kn{Both suprema exist \emph{in $\Q$}.
        Part (a) fails not because $\Q$ is too small in general but because of
        where that particular gap sits. Whether $\sup E$ lies in $E$ is a
        separate question from whether $\sup E$ exists, and running the two
        together is the most common first-week error.}
  \item Let $E$ consist of all numbers $1/n$, where $n = 1,2,3,\dots$. Then
        $\sup E = 1$, which is in $E$, and $\inf E = 0$, which is not in $E$.
\end{enumerate}

\ritem{1.10}{Definition}
An ordered set $S$ is said to have the \emph{least-upper-bound property} if the
following is true:

If $E \subset S$, $E$ is not empty, and $E$ is bounded above, then $\sup E$
exists in $S$.\kw{Both hypotheses are load-bearing, and dropping either breaks
the definition rather than merely weakening it. In $\R$: every real is an upper
bound of $\emptyset$, so $\emptyset$ has no \emph{least} one; and $\Z$ is
unbounded above, so it has none either. Any proof that produces a supremum owes
you a check of both.}

Example 1.9(a) shows that $\Q$ does not have the least-upper-bound property.

\begin{deeper}[Rudin's choice of axiom is one of several]
Completeness can be installed in an ordered field in several equivalent ways,
and which one a book takes as primitive shapes everything after it.

\smallskip
\textbf{Rudin} takes the least-upper-bound property as the defining property
(1.10), asserts a field with it (1.19), and \emph{derives} the archimedean
property from it (1.20a).

\smallskip
\textbf{Yu Pin} (Tsinghua) axiomatises $\R$ by field, order, archimedean, and
the \emph{nested interval} axiom --- every decreasing sequence of closed bounded
intervals has a common point. He observes that $\Q$ satisfies the first three
and fails only the fourth, then derives the supremum principle from it.
Archimedes must be assumed separately, because nested intervals alone does not
imply it; he proves the converse too, so over an archimedean ordered field the
two are equivalent.

\smallskip
\textbf{Tao} (\emph{Analysis I}) assumes nothing, builds $\R$ from $\Q$ as
equivalence classes of Cauchy sequences, and \emph{proves} the
least-upper-bound property as a theorem.

\smallskip
\textbf{The ECNU text} lays out the whole menu. Its \S7.1 collects six
statements --- the supremum principle, monotone convergence, nested intervals,
Heine--Borel, Bolzano--Weierstrass, and the Cauchy criterion --- and proves the
cycle $1\Rightarrow2\Rightarrow\cdots\Rightarrow6\Rightarrow1$, so that any one
of them may serve as the axiom and the other five follow.

\smallskip
That last point is worth carrying into Chapter 2. Rudin does not present these
as six faces of one property; he proves them separately and far apart --- 2.41
and 2.42 for compactness and limit points in $\R^k$, 3.11 for Cauchy sequences,
3.14 for monotone ones. Each is a theorem earned from 1.19. Knowing they are
interderivable tells you that those later proofs are not accumulating new
strength about $\R$: they are unpacking, in different vocabularies, the single
property assumed here.
\end{deeper}

We shall now show that there is a close relation between greatest lower bounds
and least upper bounds, and that every ordered set with the least-upper-bound
property also has the greatest-lower-bound property.

\begin{unpack}[The picture behind Theorem 1.11]
We are handed a hypothesis about suprema and asked to produce an infimum, so
the only move available is to find a set whose \emph{supremum} is the infimum we
want. The set of lower bounds $L$ is that set --- it sits entirely below $B$,
every element of $B$ bounds it from above, and the two meet at a single point.

\medskip
\begin{center}
\begin{tikzpicture}[x=1mm, y=1mm]
  % the set B, well clear above the axis
  \draw[seg] (49,11) -- (88,11);
  \node[lbl, anchor=south] at (68,13) {the set $B$};
  % the axis, carrying L
  \draw[axis] (-1,0) -- (94,0);
  \draw[seg] (4,0) -- (48,0);
  \node[ptfill] at (48,0) {};
  \node[lbl, anchor=north] at (24,-2) {$L$ = all lower bounds of $B$};
  % the shared boundary
  \draw[guide] (48,-6) -- (48,18);
  \node[lbl, anchor=south] at (48,18) {$\alpha=\sup L=\inf B$};
  % B presses down on L
  \draw[thin, -{Stealth[length=3.5pt]}] (62,9.5) -- (56,2);
  \node[lbl, anchor=west] at (63,6) {each $x\in B$};
  \node[lbl, anchor=west] at (63,2.6) {bounds $L$ above};
\end{tikzpicture}
\end{center}
\figcap{The two sets share one boundary point, and the theorem says that point
is both $\sup L$ and $\inf B$.}
\end{unpack}

\ritem{1.11}{Theorem}
\emph{Suppose $S$ is an ordered set with the least-upper-bound property,
$B \subset S$, $B$ is not empty, and $B$ is bounded below. Let $L$ be the set of
all lower bounds of $B$. Then}
\[ \alpha = \sup L \]
\emph{exists in $S$ and $\alpha = \inf B$.}

\emph{In particular, $\inf B$ exists in $S$.}

\begin{proof}
Since $B$ is bounded below, $L$ is not empty. Since $L$ consists of exactly
those $y \in S$ which satisfy the inequality $y \le x$ for every $x \in B$, we
see that every $x \in B$ is an upper bound of $L$.\kn{This single sentence is
the whole idea. We are handed a hypothesis about suprema and asked for an
infimum, so we must find a set whose \emph{supremum} is the infimum we want.
$L$ is that set: $B$ presses down on it from above, and the two meet at one
point.} Thus $L$ is bounded above. Our hypothesis about $S$ implies therefore
that $L$ has a supremum in $S$; call it $\alpha$.

If $\gamma < \alpha$ then (see Definition 1.8) $\gamma$ is not an upper bound of
$L$, hence $\gamma \notin B$.\kn{Compressed to the point of opacity. Unpacked:
$\gamma$ failing to be an upper bound of $L$ produces some $y \in L$ with
$y > \gamma$. Since $y$ is a lower bound of $B$, every $x \in B$ satisfies
$x \ge y > \gamma$. So $\gamma$ itself cannot belong to $B$.} It follows that
$\alpha \le x$ for every $x \in B$. Thus $\alpha \in L$.

If $\alpha < \beta$ then $\beta \notin L$, since $\alpha$ is an upper bound of
$L$.

We have shown that $\alpha \in L$ but $\beta \notin L$ if $\beta > \alpha$. In
other words, $\alpha$ is a lower bound of $B$, but $\beta$ is not if
$\beta > \alpha$. This means that $\alpha = \inf B$.
\end{proof}

\begin{proofwalk}[How the proof flows]
  \item \textbf{Choose the method.} The hypothesis speaks of suprema, the goal
        of an infimum, and the least-upper-bound property is the only
        existence tool in the chapter. So the method is forced --- pattern 3
        of the strategy index: \emph{find a set whose supremum is the number
        you want}, then take its sup. The candidate is $L$, the set of all
        lower bounds of $B$; the picture before the theorem shows why it is
        right: $L$ ends exactly where $B$ begins.
  \item \textbf{Buy the supremum.} The lub property sells only to sets that
        are nonempty and bounded above, so both tickets must be produced.
        $L \ne \emptyset$ is the hypothesis ``$B$ is bounded below''
        restated. $L$ bounded above is the proof's one genuine observation:
        every $x \in B$ bounds $L$, by the very definition of ``lower
        bound.'' Payment accepted: $\alpha = \sup L$ exists.
  \item \textbf{Show $\alpha$ is a lower bound} (first clause of the goal).
        Rudin argues by testing every $\gamma < \alpha$: clause (ii) of 1.8
        says $\gamma$ is not an upper bound of $L$, so some $y \in L$ exceeds
        $\gamma$; since $y$ underlies all of $B$, so does $\gamma$, strictly
        --- hence $\gamma \notin B$. Nothing in $B$ lies below $\alpha$, which
        is precisely $\alpha \in L$.
  \item \textbf{Show nothing larger works} (second clause). Any
        $\beta > \alpha$ exceeds the supremum of $L$, hence $\beta \notin L$:
        it fails to be a lower bound. This step uses clause (i) of 1.8 ---
        $\alpha$ \emph{is} an upper bound of $L$.
  \item \textbf{Read off the conclusion.} Steps 3 and 4 are exactly the two
        clauses of Definition 1.8, mirrored for the infimum, so
        $\alpha = \inf B$. Note the audit: clause (ii) of the supremum was
        spent in step 3, clause (i) in step 4 --- each exactly once, the
        usual sign that $L$ was the right set to build (compare the same
        remark at 1.21).
\end{proofwalk}

\begin{center}
\begin{tikzpicture}[x=1mm, y=1mm]
  \node[rmbox, text width=27mm] (h) at (0,42)
    {$B \ne \emptyset$,\\bounded below\\ \textit{(hypothesis)}};
  \node[rmbox, text width=27mm] (l) at (60,42)
    {$L \ne \emptyset$, bounded\\above by every $x \in B$};
  \draw[rmarrow] (h) -- node[lbl, anchor=south, inner sep=2.5pt]
    {definition of $L$} (l);
  \node[rmbox, text width=24mm] (s) at (30,22) {$\alpha = \sup L$ exists};
  \draw[rmarrow] (l) -- node[lbl, anchor=west, inner sep=3pt]
    {\lub{} property} (s);
  \node[rmbox, text width=32mm] (b1) at (0,0)
    {$\gamma < \alpha \Rightarrow \gamma \notin B$:\\
     $\alpha$ is a lower bound};
  \node[rmbox, text width=32mm] (b2) at (60,0)
    {$\beta > \alpha \Rightarrow \beta \notin L$:\\
     no larger one is};
  \draw[rmarrow] (s) -- node[lbl, anchor=east, inner sep=3pt]
    {1.8(ii)} (b1);
  \draw[rmarrow] (s) -- node[lbl, anchor=west, inner sep=3pt]
    {1.8(i)} (b2);
  \node[rmbox, text width=22mm] (c) at (30,-21) {$\alpha = \inf B$};
  \draw[rmarrow] (b1) -- (c);
  \draw[rmarrow] (b2) -- (c);
\end{tikzpicture}
\end{center}
\figcap{The proof as a flow diagram: hypotheses qualify $L$, the
least-upper-bound property mints $\alpha$, and the two clauses of Definition
1.8 each certify one half of ``$\alpha = \inf B$.''}

\begin{deeper}[Why this is not a triviality]
The argument runs equally well in reverse, so for ordered sets the
least-upper-bound property and the greatest-lower-bound property are
\emph{equivalent}. That is why Rudin never states a separate greatest-lower-bound
axiom for $\R$: Theorem 1.19 gives him one and Theorem 1.11 hands him the other.

Every later appearance of $\inf$ rests on this --- lower Riemann--Stieltjes sums
in 6.2, $\liminf$ in 3.16, the distance from a point to a set in the exercises.
It is a two-paragraph theorem that the rest of the book quietly leans on.
\end{deeper}
